% ── SECTION 6: DISCUSSION (~600 words) ───────────────────────────────────────

\section{Discussion}
\label{sec:disc}

\subsection{Physical Interpretation of the Affinity Tensor}

The Affinity Tensor~$\bar\alpha$ provides an operationalized measure of
relational coherence that is distinct from existing entanglement measures
in both its definition and its purpose.
Concurrence~\cite{Wootters1998} and von Neumann entropy~\cite{vonNeumann1932}
are static measures: they quantify entanglement at a snapshot in time.
$\bar\alpha(t)$ is a dynamical measure: it tracks the evolution of
off-diagonal coherence under environmental coupling, decaying from~1
to~0 as the system transitions from quantum to classical behavior.

For two-qubit pure states, Proposition~\ref{prop:correspondence} establishes
$\bar\alpha = \mathcal{C}$, so the Affinity Tensor and concurrence agree.
The operational difference appears under noise: $\bar\alpha(t)$ provides a
single scalar trajectory that directly quantifies the relational coherence
budget available to a gate at each moment of the circuit's execution, whereas
concurrence characterizes static entanglement of a state.
This makes $\bar\alpha$ the natural control parameter for hardware-adaptive
entangling gates, which is precisely the role $\URA(\alpha)$ fills.

\subsection{The Gate-Time Advantage and its Interpretation}

The simulation results of Section~\ref{sec:sim} confirm a direct prediction
of the \RCT: for tasks requiring partial entanglement ($\bar\alpha < 1$),
operating at reduced~$\alpha$ preserves relational coherence that a fixed
maximally entangling gate would destroy.
The fidelity advantage~---~$\Delta F \in [+0.019, +0.040]$ at
$p_{\rm base} = 0.10$ across all noise models~---~is not a fine-tuning
artifact; it arises from the fundamental scaling of decoherence with
gate duration (Eq.~\ref{eq:peff}) combined with the continuous
parameterization of $\URA$ through the Weyl chamber.

This result has a clear interpretation in terms of the four propositions.
Conservation of Relational Potential (Proposition~\ref{prop:conservation})
guarantees that the coherence ``spent'' during a gate operation transfers
to system-environment correlations.
A gate that operates below maximal entanglement capability spends less
of this coherence budget per operation, leaving more available for
subsequent circuit layers.
The Affinity Tensor~$\bar\alpha$ is therefore not only a diagnostic
measure but a \emph{resource}: controlling it directly through~$\alpha$
is equivalent to controlling the coherence expenditure of the gate.

\subsection{Scope and Limitations}

The \RCT~formalizes the relational structure of quantum mechanics established
by Bell's theorem and relational quantum mechanics; it does not resolve the
quantum measurement problem.
It provides a precise formal framework within which the measurement problem
can be stated~---~a definite outcome is the establishment of a relational fact
(Proposition~\ref{prop:determination})~---~but the question of \emph{which}
outcome is selected in a given run of an experiment remains outside the
framework's scope, as it does for all no-collapse interpretations.

The simulation results employ a gate-time-proportional noise model appropriate
for Heisenberg exchange interactions on superconducting hardware~\cite{Krantz2019}.
On hardware where gate duration does not scale with~$\alpha$ (e.g., where
$\URA$ must be decomposed into fixed native gates), the fidelity advantage
depends on the depth of the decomposition and may be reduced.
Direct hardware calibration of the Heisenberg exchange parameter, as
demonstrated for \textsc{iswap}-family gates, would realize the full advantage.

\subsection{Future Directions}

Three extensions are immediate.
First, experimental implementation of~$\URA$ on superconducting qubit
platforms~\cite{Krantz2019} via direct calibration of the Heisenberg
exchange coupling, bypassing gate decomposition overhead.
Second, extension to multi-qubit relational networks: generalizing the
Affinity Tensor to $n$-qubit systems $\alpha(S, \{E_i\})$ and deriving the
corresponding $n$-qubit Relational Affinity unitaries parameterized by the
full tensor.
Third, investigation of topological quantum error correction codes
parameterized by~$\bar\alpha$: the Affinity Tensor may provide a natural
measure for the coherence threshold in surface codes.
The philosophical implications of the \RCT~---~including the argument that
process ontology is formally required (not merely permitted) by quantum
mechanics~---~are developed in the companion paper~\cite{Adams2026Philosophy}.
